\section{A smaller conductor constant and a shifted detector coordinate}

This note starts from the immutable fourth-round package with builder hash
\texttt{ca171bf25610cb7da003ed58b4c4d43799a67fd4b5d6ca1b8fbac74a2dcc7e0e}.
The upstream manuscript is pinned at commit
\texttt{036cebe4f1a4aa7ca89820a9b48f3629b4c364d6}.
All global geometric, modularity, height-transfer, group, simultaneous-prime,
pole and quaternion ingredients remain inherited. This note checks changed
applications of their actual proofs; computational checks of its constants
are not substitutes for those ingredients.

\subsection{Use the full three-torsion already defined over the fiber field}

The actual geometric construction in
\texttt{manuscripts/quartic\_v10.tex}, lines 6933--7006, provides a
normal field $L/\mathbb Q$ with $[L:\mathbb Q]\mid576$, a maximal-order
quaternion action on $A/L$, and full three-torsion markings defined over
$L$. In particular all of $A[3]$ is $L$-rational.

Raynaud's criterion, in Silverberg--Zarhin, \emph{Semistable reduction
and torsion subgroups of abelian varieties}, Ann. Inst. Fourier 45 (1995),
Theorem 3.5 on printed page 406, applies to full $n$-torsion for $n\ge3$
prime to the residue characteristic. The primary published page was
checked directly, including the inequality $n\ge3$.
The source already invokes this theorem with three-torsion in lines
883--885 and with eleven-torsion in lines 7153--7156.
Thus $A$ is semistable over $L$ at $2,5,7$.

The source's quaternionic torus argument, lines 7105--7126, upgrades
this to good reduction: a positive-rank torus in the semistable special
fiber would give a unital embedding of the quaternion division algebra
$B$ into $M_t(\mathbb Q)$ for $1\le t\le2$, which is impossible.
At $3$, retain the previous field $E=L(A[11])$. The fourth-round
matrix-unit argument gives
\[
 [E:\mathbb Q]\mid2^{10}3^3 5^2 11,
\]
and eleven-torsion gives good reduction there.

Let $\chi,b$ be precisely the scalar character and cochain constructed
in the source's (LR2). On $G_L$, $b$ is a conjugation-invariant character
whose square is $\chi$. Put $F_0=L\mathbb Q(\chi)$ and let $F'_0$
kill the quadratic restriction of $b$ to $G_{F_0}$. This is a normal
extension, by the same argument used in the source for $E\mathbb Q(\chi)$.
Over $F'_0$, the scalar correction is trivial and the compatible
representation is the Tate representation of $A$.

At $2$, the two scalar extensions each add ramification index at most
two; the local valuation bound is therefore $v_2(e_{F'_0})\le6+2=8$.
At $5,7$ both added extensions are tame and $[L:\mathbb Q]$ is prime
to those residue characteristics. At $3$, use the old $E$-based extension;
its two added extensions are tame, so $v_3(e)\le3$. It is legitimate to
use different auxiliary fields at different primes to bound one local
conductor. The resulting wild-index bounds are
\[
 (k_2,k_3,k_5,k_7)=(8,3,0,0).
\]

\subsection{Retain the sharp term in the different formula}

For the finite Galois local extensions just used, the source proves
\[
 \delta\le e-1+e\,v_p(e),\qquad
 \delta/e=1-1/e+\int_0^\infty(1-|G^u|^{-1})\,du.
\]
These are the conventions in lines 7367--7386. Subtracting the common
$1-1/e$ term gives an integral at most $k=v_p(e)$.
For $u>0$, every nontrivial $G^u$ is a $p$-group, so the integrand is
at least $(p-1)/p$ up to the last positive upper break. The value at
$u=0$ and the choice of endpoint convention have measure zero and do not
change this estimate. The last break is at most $pk/(p-1)$.

The two-dimensional representation is unramified over the corresponding
auxiliary field. Its Swan conductor is therefore at most $2pk/(p-1)$,
and its tame contribution is at most two. If $k=0$, it is tame. Hence
\[
 a_2\le34,\qquad a_3\le11,\qquad a_5\le2,\qquad a_7\le2.
\]
Use the source's companion representation at three to compute the
conductor at five, exactly as before. At every $p>7$ the original bound
$a_p\le2$ is unchanged, with zero conductor outside its support $T$.
Since $R$ contains $2,3,5,7$, the already present factor $R^2$ absorbs
two powers of each. Consequently
\[
 \boxed{N\le C R^2,\qquad C=2^{32}3^9=84537841287168<2^{47}.}
\]
This refinement changes neither the modular representation nor its
primitive-level interpretation.

\subsection{Choose the detector coordinate and guard together}

Replace the old fractional-linear coordinate by
\[
 \phi(t)=-49+1/t,
\]
which sends the five new contact values
\[
 \Theta=\{0,1/49,1/50,1/2450,32/1225\}
\]
to $\{\infty,0,1,2401,-343/32\}$. Its integral matrix has determinant
$-1$. All contact denominators and nonzero pairwise differences have
prime support in $\{2,3,5,7\}$.

For each chosen positive-square scale $A$, put
\[
 t=A\lambda,\qquad\delta=98t-1,\qquad
 F(t)=\prod_{\theta\in\Theta}(t-\theta),\qquad
 e=(A\delta)^5/F(t),\qquad R=1/(e\delta).
\]
The values of $\delta$ on the five contacts, in the displayed order, are
\[
 -1,\quad1,\quad24/25,\quad-24/25,\quad39/25.
\]
Their only extra prime is $13$. This agrees with the exception already
needed by the retained input logarithm.

The positivity chart becomes
\[
 \lambda=32/1225+u^2+v^2.
\]
The largest contact is $32/1225$. All three scales are at least one,
so $t$ lies at or above this value, where $\delta\ge39/25>0$.
The defining equation $eF=(A\delta)^5$ excludes a contact value and
therefore gives $F>0$, $e>0$. For every positive order $m$,
\[
 q_m=e^{2m-1}/\delta^3>0.
\]
Every real detector factor $1+q_mD_j^2$ is at least one. Thus all prior
nonzero trace, norm and pole-input arguments remain valid; no new guard
is silently assumed.

At every prime outside $\{2,3,5,7,13\}$, a prime in the transformed
contact radical has finite contact with exactly one element of $\Theta$,
and $\delta$ is a unit there. Contact at infinity for $t$ does not occur:
$\phi(\infty)=-49$ is distinct from all five target contacts modulo these
primes. The matrix is unimodular, so the same odd-contact radical is
obtained after applying $\phi$. The usual one-place argument applies
because
\[
 1+R\left(\frac{D/\delta^{m+1}}{R^m}\right)^2
       =1+\frac{e^{2m-1}}{\delta^3}D^2.
\]
The identity and its valuation proof do not require the old guard $t-1$.

\subsection{Check the actual logarithm orders}

Use $G_{12}$ for the shared alternate denominator and for the three
Lucas-label numerators, with high detector order $44$. Retain $G_2$
for the input numerator and low cross order $12$.
The exact circuit identities and global nonvanishing of these polynomial
approximants are proved in the separate geometry note.

Set
\[
 \mathcal S=\{2,3,5,7,13,17\},\qquad B_{\mathcal S}=46410.
\]
All primes outside this set are at least eleven. For a formal point
of depth $r\ge1$, the source identity gives $\log_E=zG(w)$ with
$v(w)=4r$ and $v(a_j)\ge-(4j+1)/(p-1)$.
For $G_{12}$, every omitted $j\ge13$ satisfies
\[
 (4j+1)/10\le4(j-11),
\]
so its relative error has valuation at least $44r$.
For $G_2$, the first omitted coefficient is
$a_3=-14/195$, integral at all consumed primes. For $j\ge4$,
$(4j+1)/10\le4(j-3)$, giving relative precision $12r$.
For raw $z$, $a_1=-2/5$ is integral and the same estimate for $j\ge2$
gives $\log_E/z\in1+p^{4r}\mathcal O$.
These tail bounds tend to infinity and control the entire series.
They hold at split and inert unramified places. This explicitly checks
the retained input logarithm; omitting $13$ would fail its order-twelve
consumer.

The primary root test at high order $44$ gives primary formal depth
at least $44$. Weak logarithm-unit estimates and the fact that one of
$h,h-1$ is a $p$-unit transfer integrality before relative errors are
made absolute. The three Lucas labels are then compared at order $44$,
and the input label at order $12$. The primary error is at least $176$.
The existing identity-input branch uses its own cross-test order and
does not form $\log(O)/z(O)$.

\subsection{Height constants, scales and the anchor}

The finite-support geometry, (LR1)--(LR5), CM branch and exponent-nineteen
height proof are the actual ones audited in the fourth-round note.
They apply with the smaller support and the new conductor constant:
\[
 H=2\cdot2^{109}C^6B_{\mathcal S}^{19}+128,
 \qquad h(x)\le H M_{\mathcal S}(x)^{19}.
\]
The exact $H$ has bit length $683$. The elementary bound used below is
\[
 H<2^{697},\qquad697=2+109+6\cdot47+19\cdot16.
\]
This still dominates the inherited CM estimate.

For the selected one-coordinate scale reduction, retain $1,C_0^2$ and
replace $(1+C_1^2)^2$ by $(C_0+C_1)^2$, where
$C=C_0+C_1\rho$ is the product of the four root abscissas.
At both real embeddings $C>d^4$, and $d^4=193-132\rho$.
The positive weights $(1\pm1/\sqrt2)/2$ therefore give $C_0+C_1>61$.
All scales remain greater than or equal to one.

The coordinate change satisfies $h(t)\le h(\phi(t))+\log50$.
If all three radicals divide a nonzero integer $N$ to order $k$, put
$T=(H+\log50)|N|^{19/k}$. Then $h(C_0),h(C_0+C_1)\le T$ and
\[
 h(C_1)\le2T+\log2,\qquad h(C)\le3T+\tfrac52\log2.
\]
Since $H\ge4$ and $3\log50+\tfrac52\log2<16$, the old constant
$7H$ still suffices. The product-height identity and root lower bound
therefore retain the exact consumer
\[
 B^2-6<7H|N|^{19/k},\qquad B=\max(|n|,|a|,|t|,|b|).
\]

Use scalar order $m=28$. Starting from $q_1=r^{-1}$, four coordinates
give
\[
 rq_3=q_1^2,\qquad rq_7=q_3^2,\qquad
 q_{14}=q_7^2,\qquad q_{28}=q_{14}^2.
\]
The previously reviewed scalar norm recursion is uniform in positive
orders, so $M_{\mathcal S}(\beta)^{28}\mid a-1$ and
$h(\beta)\le H|a-1|^{19/28}$. The separately reviewed unnormalized
scalar proxy retains exactly this index valuation conclusion, including
the prime $41$; it needs no additional exceptional prime.

Set the anchor coefficient to $S_0=3^{2^{13}}$, with
$C_a=3^{2^{12}}$. The unchanged norm-anchor calculation gives
\[
 0\ne a-1\in2\cdot3^{2^{11}-1}\mathbb Z.
\]
Write $c=|a-1|$. Then
\[
 \log_2c>3071.5,\qquad\log_2|a|>3070.5.
\]
The scalar margin follows from
\[
 (9/28)3071.5>987>704=7+697,
\]
which gives $h(\beta)<c/80\le|a|/40$.
Also $\log(40H)<703<|a|/40$; the last inequality follows already
from $40\cdot703<2^{15}<|a|$.

For the Lucas polynomial $F$, if $F\ne0$, its order-$44$ divisibility
and $|F|<5B^4$ give
\[
 B^{3/11}<7H5^{19/44}+6<24H.
\]
But
\[
 (3/11)3070.5>837>702=5+697>\log_2(24H),
\]
contradicting $B\ge|a|$. Hence $F=0$, and the inherited Lucas theorem
gives $\log B>|a|-3$. For $n\ne\beta$, the order-twelve input comparison
gives the unchanged contradictory upper bound
\[
 \log B<\tfrac{12}{5}\log(40H)+\tfrac{19}{5}h(\beta)
 <\tfrac{31}{200}|a|<|a|-3.
\]
Thus all soundness consumers are satisfied without importing the former
tenth-power margin as an unnecessary common requirement.

\subsection{Completeness and the new fixed constants}

Keep the source's order of choices: scalar preparation, roots and positive
square scales, continuity modulus, then contact primes. The old order-$32$
scalar preparation supplies order $28$. Choose contact primes greater than
$76176$ with the new fixed contact positions and the same Gaussian classes.
All new exclusions and differences are fixed before prime selection.
The separation proof still gives the required coprimalities, so the primary
progression uses $\Pi^{43}$ rather than $\Pi^{39}$, yielding depth at
least $44$. Opposite points retain depth one.

The positivity chart has canonical witnesses: $32/1225$ is a contact of
the unit-scale detector. The selected contact formula gives
$\lambda-32/1225=p/v$, where $v=D_0^2w^2$ is a square and the selected
$p$ is a Gaussian norm. Thus this difference is a sum of two rational
squares. Arbitrary positivity alone would not establish this assertion.

The product formula for $R$ still has square denominator $v^5\delta^6$,
and hence the same odd valuations and local squareclasses. The successful
place makes the reciprocal trace a unit at selected contact primes, exactly
as before. Real asymptotics now read
\[
 R\sim98^{-6}A^{-6}\lambda^{-1},\quad W=1+O(\lambda^{-3}),
 \quad B_{\rm pole}\sim(4\cdot98^6A^6\lambda)^{-1}.
\]
At $p=2,3,17$, let $\nu=v_p(v)$, $\eta=v_p(A)$ and
$\alpha=v_p(98)$. For sufficiently large prescribed $\nu$,
\[
 v_p(\delta)=\alpha+\eta-\nu,\qquad
 v_p(R)=\nu-6\eta-6\alpha,
\]
and the nonconstant detector factors have valuation at least
\[
 3\nu+(10m-8)(\eta+\alpha)-2c_p.
\]
This tends to infinity for every retained fixed order. Accordingly the
source's small-prime pole bounds and real upper bound $B_{\rm pole}<17$
hold after increasing only finite preparatory constants. No new
simultaneous-prime hypothesis or circular choice is introduced.

The standard constant chain now runs from exponent ten through thirteen,
using four coordinates. An intermediate chain can use twelve through
thirteen, using two coordinates. These count statements require the
ordinary quadratic lifting audit; the fixed-coefficient core has no extra
coordinate for the changed numerical anchor.
